\documentclass[11pt]{article}

\usepackage[margin=1in]{geometry}
\usepackage[T1]{fontenc}

\title{Frontal Cortex Dynamics Suggest Temporal Misalignment Between Sensory Evidence and Choice}

\date{}

\begin{document}

\maketitle
\thispagestyle{empty}

\begin{abstract}
Neural ensembles in mouse frontal cortex are essential for maintaining choices in short-term memory (Guo et al.\ 2014). Modeling work explains choice maintenance through attractor dynamics, where stimulus onset triggers a landscape that stabilizes choice representations until response (Inagaki et al.\ 2019). However, this model assumes tight temporal alignment between stimulus presentation and decision, an assumption that may fail when the stimulus timing is unpredictable. To explore this, we developed a three-choice visuospatial delayed-response task for freely moving mice with variable stimulus and delay durations. Choice accuracy increased with stimulus duration and decreased with delay length, suggesting both perceptual and forgetting errors. Importantly, mice still responded on trials where the stimulus was never presented, indicating that decisions can be initiated independently of sensory input. Motivated by this observation, we extended the two-choice attractor network model to a three-choice circuit with three excitatory populations coupled via global inhibition. The network received two inputs: a transient stimulus-selective signal and a non-selective ramping input representing an internal urgency to decide. These models showed that task performance depended on the temporal alignment between the sensory input and urgency buildup. Misalignment produced a previously unconsidered error occurring when urgency peaked after stimulus offset, resulting in delayed choice encoding. Consistent with this prediction, population activity recordings in the secondary motor cortex revealed that choice encoding emerged later in error trials. Together, these results show that decisions in this task are not necessarily aligned with the stimulus presentation, and the relative timing between stimulus and choice determines task performance.
\end{abstract}

\end{document}